\documentclass[11pt]{article}
\usepackage{amsmath}
\usepackage[hidelinks]{hyperref}

\begin{document}

\section{Methodology}
\label{sec:sdpo-methodology}

We adapt Self-Distillation Policy Optimization (SDPO)
\cite{hubotter2026sdpo} to multilingual mathematical reasoning.
During training, a verifier identifies successful attempts, and reasoning
from verified English attempts provides additional context to a self-teacher.
The student learns from the teacher's next-token distributions on failed
attempts and from a group-relative policy objective on successful attempts.
At inference, the student receives the problem in the requested language
without a reference or verifier feedback.

\paragraph{Parallel problems and outcome verification.}
We organize mAceReason-Math examples into parallel groups, each containing
one version of the same problem per training language. A training batch
contains two complete groups. For each localized prompt \(x_i\), the
rollout policy samples 32 responses. The verifier extracts the last boxed
answer from the final response and assigns a binary reward. It first checks
for an exact match and then checks mathematical equivalence if needed.

\paragraph{English reference selection and teacher conditioning.}
For each failed response, we sample uniformly from eligible rollouts
generated from the English prompt in the same parallel group. A candidate
must have a correct final answer, a nonempty closed reasoning trace, and
no truncation or repeated-final-answer flag. We discard the final response
following the reasoning trace and cut the trace before the line containing
the first detected explicit answer cue. We append the extracted reasoning
to the original localized prompt and ask the teacher to solve the problem
using a similar approach in the target language. If no eligible reference
is available, the teacher can instead receive verifier feedback about
truncation or a missing boxed answer. Failed responses with neither form
of context, or whose augmented teacher prompts exceed the token budget,
are masked out of the distillation objective.

\paragraph{Token-level self-distillation.}
Let \(x_i\) be the original math problem prompt in the target language
\(\ell_i\). Let \(y_i=(y_{i,1},\ldots,y_{i,T_i})\) be a response sampled
from the rollout policy conditioned on \(x_i\), where \(T_i\) is the number
of response tokens and \(y_{i,<t}=(y_{i,1},\ldots,y_{i,t-1})\) is the
response prefix before position \(t\).

For a failed response \(y_i\) with an eligible English reference, let
\(j(i)\) index the selected verified English rollout from the same
parallel group. Its extracted
reasoning is \(r_i^{\mathrm{en}} =
\operatorname{ExtractReasoning}(y_{j(i)})\). The augmented teacher prompt is
\begin{equation}
\hat{x}_i =
\underbrace{x_i}_{\text{original localized prompt}}
\;\Vert\;
\underbrace{\operatorname{ReferenceSection}
  (r_i^{\mathrm{en}},\ell_i)}_{\text{English reasoning and its framing}}
\;\Vert\;
\underbrace{\operatorname{Reminder}(\ell_i)}_{\text{target-language reminder}}.
\label{eq:sdpo-teacher-prompt}
\end{equation}
Here \(\Vert\) denotes concatenation within the user message. The reference
section contains the English reasoning trace and instructs the teacher to
apply a similar approach in \(\ell_i\). For a feedback-only failure,
\(\hat{x}_i\) is built from \(x_i\) and the verifier feedback using the
feedback prompt template.

Using the teacher prompt from Eq.~\eqref{eq:sdpo-teacher-prompt}, at each
valid response position \(t\), the student and teacher predict over
the same vocabulary \(\mathcal V\) using the same previously sampled tokens
\(y_{i,<t}\). For any candidate next token \(v\in\mathcal V\),
\begin{align}
p_{i,t}(v) &= \pi_\theta(v\mid x_i,y_{i,<t}), \\
q_{i,t}(v) &= \pi_{\bar\theta}(v\mid\hat{x}_i,y_{i,<t}).
\label{eq:sdpo-next-token-distributions}
\end{align}
The English reference thus enters through the teacher's conditioning context
at every response position. The teacher forward pass is detached, and its
parameters track the student through an exponential moving average,
\(\bar\theta\leftarrow 0.95\bar\theta + 0.05\theta\).

Let \(S_{i,t}\) contain the student's 512 highest-probability tokens at
position \(t\). For each distribution \(u\in\{p_{i,t},q_{i,t}\}\), we
retain \(u(v)\) for \(v\in S_{i,t}\) and assign the remaining mass
\(1-\sum_{v\in S_{i,t}}u(v)\) to a single tail category. Denote the resulting
distributions by \(p^*_{i,t}\) and \(q^*_{i,t}\). The per-token divergence is
\begin{equation}
d_{i,t}=\operatorname{JSD}_{\alpha}(p^*_{i,t},q^*_{i,t}).
\label{eq:sdpo-token-divergence}
\end{equation}
The generalized Jensen--Shannon divergence uses weights \(1-\alpha\) for
the student and \(\alpha\) for the teacher, with both KL terms evaluated
against their weighted mixture. The default \(\alpha=0.5\)
yields the symmetric Jensen--Shannon divergence.

Let \(b_{i,t}\) indicate a valid response token and \(s_i\) indicate a
failed response with usable teacher context. Under the hybrid policy update,
the distillation term is
\begin{equation}
\mathcal L_{\mathrm{SDPO}} =
\frac{\sum_{i,t} b_{i,t}s_i\,w_{i,t}\,d_{i,t}}
     {\max(1,\sum_{i,t} b_{i,t})}.
\label{eq:sdpo-loss}
\end{equation}
The weight \(w_{i,t}\) includes the detached current-to-old policy
probability ratio for \(y_{i,t}\), capped at 2, and multiplies it by the
token-level rollout importance weight when available.

\paragraph{Hybrid policy update.}
For each localized prompt, we compute a GRPO advantage by subtracting the
mean reward across its sampled responses from each response's reward.
The mean includes both successful and failed attempts, and the advantage
is used without standard-deviation normalization. Successful responses
contribute a clipped policy-gradient loss. Failed responses with usable
teacher context contribute the distillation loss. Both contributions are
normalized by the full count of valid response tokens in the loss
reduction. The actor also receives a KL penalty against the fixed
reference policy with coefficient \(10^{-3}\).

\paragraph{Training configuration.}
The default mAceReason recipe uses Nemotron-Research-Reasoning-Qwen-1.5B,
up to 500 parallel problem groups, and one epoch at a learning rate of
\(3\times10^{-6}\). The default languages are Bengali, English, Spanish,
Telugu, and Thai. The augmented teacher prompt has a budget of 10,752
tokens before the sampled response is appended.

\begin{thebibliography}{1}
\bibitem{hubotter2026sdpo}
J.~H\"ubotter et al.
\newblock Reinforcement Learning via Self-Distillation.
\newblock \href{https://arxiv.org/abs/2601.20802}{arXiv:2601.20802}, 2026.
\end{thebibliography}

\end{document}
